\section{JSON Schema}

Source: \code{schemas/difftracker.schema.json}.

Lean mirror: \code{Defns.lean} (structured diff types), \code{JsonParse.lean} (parser). Lean declarations use the \code{VersionControl} namespace.

The schema defines the wire format for one \code{Diff} (Definition~\ref{def:diff}). Commit ancestry is stored separately.

All terms are defined in Section 1. This section maps schema keywords to constraints.

\subsection{Top-Level Object}

The top-level object must contain exactly these fields (\code{"required"} and \code{"additionalProperties": false}):

\begin{itemize}
  \item \code{agent}: \code{integer}, \code{minimum: 1}.
  \item \code{file}: \code{string}, \code{minLength: 1}.
  \item \code{timestamp}: \code{string} with \code{format: date-time}.
  \item \code{base_line_count}: \code{integer}, \code{minimum: 0}.
  \item \code{diffs}: \code{array} of objects matching \code{"\#/\$defs/hunk"}.
\end{itemize}

\subsection{Hunk Object}

Each element of \code{diffs} is a hunk (\code{"\$ref": "\#/\$defs/hunk"}). It must contain exactly these fields (\code{"required"} and \code{"additionalProperties": false}):

\begin{itemize}
  \item \code{seq}: \code{integer}, \code{minimum: 1}.
  \item \code{op}: one of \code{"insert"}, \code{"delete"}, \code{"replace"} (\code{"enum"}).
  \item \code{range}: array of exactly two integers $\geq 1$ (\code{minItems: 2}, \code{maxItems: 2}, \code{items: \{type: integer, minimum: 1\}}).
  \item \code{before}: array of strings.
  \item \code{after}: array of strings.
\end{itemize}

The \code{allOf} array with \code{if}/\code{then} clauses enforces operation shapes:
\begin{itemize}
  \item \code{insert}: \code{before} empty, \code{after} non-empty.
  \item \code{delete}: \code{before} non-empty, \code{after} empty.
  \item \code{replace}: both \code{before} and \code{after} non-empty.
\end{itemize}

\subsection{Range Semantics}

\code{range} is a half-open interval \code{[start, stop)} with 1-based indexing.

\begin{itemize}
  \item \code{start} is included, \code{stop} is excluded.
  \item Width = \code{stop - start}.
  \item \code{start = stop} denotes an insert position.
  \item Valid for base file of \code{N} lines: \code{1 <= start <= stop <= N+1}.
\end{itemize}

Examples:
\begin{itemize}
  \item \code{[2,4]} affects lines 2 and 3.
  \item \code{[1,1]} inserts before line 1.
  \item \code{[N+1,N+1]} appends to end of file.
\end{itemize}

\subsection{Lean Constraints}

The schema enforces structure. Lean enforces additional invariants in \code{Diff.checkSchema} and \code{Hunk.checkSchema}:

\begin{itemize}
  \item \code{agent >= 1} and \code{file != ""}.
  \item No two hunks overlap within one diff, and no two inserts share the same position within one diff.
  \item \code{stop <= base_line_count + 1}.
  \item For \code{delete} and \code{replace}: \code{before.length = stop - start}.
  \item \code{start <= stop}.
\end{itemize}

\subsection{Schema vs Lean Boundary}

\begin{itemize}
  \item Schema: rejects unknown fields (\code{additionalProperties: false}), enforces types, \code{enum}, array lengths, basic shapes via \code{allOf/if/then}.
  \item Lean: stores \code{timestamp} as \code{String}, requires the five top-level fields when parsing, and validates inter-field relations and non-overlap.
\end{itemize}

JSON Schema defines the wire format. Lean defines semantic validity.

\subsection{Base Alignment}

\code{Diff.checkAgainstBase} and \code{Hunk.baseAligned} verify:

\begin{itemize}
  \item \code{before} exactly matches \code{slice(base, range)}.
  \item \code{base_line_count = base.length}.
\end{itemize}

These checks occur before \code{applyDiff}.

\subsection{Example}

Base file:
\begin{verbatim}
name = Alice
greeting = hello
print greeting
\end{verbatim}

Diff:
\begin{lstlisting}[style=jsonstyle]
{
  "agent": 1,
  "file": "hello.py",
  "timestamp": "2026-04-21T10:00:00Z",
  "base_line_count": 3,
  "diffs": [
    {
      "seq": 1,
      "op": "insert",
      "range": [1, 1],
      "before": [],
      "after": ["# my script"]
    },
    {
      "seq": 2,
      "op": "replace",
      "range": [2, 3],
      "before": ["greeting = hello"],
      "after": ["greeting = hello world"]
    }
  ]
}
\end{lstlisting}

Result after apply:
\begin{verbatim}
# my script
name = Alice
greeting = hello world
print greeting
\end{verbatim}
